% 原书分章 PDF 第 17 页：第 28.6 节，完整结论。
\section{结论}

我们从理论与模拟结果两方面，展示了靶向自助法在估计抽样分布性质方面的有效性。模拟说明了以下问题：当 TMLE 的初始估计基于超级学习器时，对该 TMLE 估计应用非参数自助法会发生什么。具体而言，非参数自助样本中的重复取值破坏了交叉验证的样本划分，使交叉验证风险估计的行为更接近回代估计。这导致超级学习器在非参数自助样本上选择过拟合模型。靶向自助法从连续分布估计中抽样，而且该分布已靶向于感兴趣的参数，因此不会产生破坏交叉验证的重复取值，从而生成表现可以接受的置信区间。我们展示了，在靶向学习情境中，靶向自助法优于非参数自助法。

还需要进一步研究靶向学习中的自助法置信区间问题。对连续结局变量进行模拟研究，尤其让其误差分布具有偏斜，会在多个方面很有意义。首先，它能验证靶向自助法用于连续结局的方法；理论告诉我们，即使误差分布设定错误，该方法也应当相合。其次，它能让我们在较小样本量下抽样分布可能更加偏斜的情形中，考察二阶项的大小。另一种推广是研究更多类型的自助法置信区间，尤其是倾斜区间与自动百分位区间（DiCiccio 和 Romano，1990）。

我们还希望强调另一种在相当一般的情形下渐近相合的自助法。这一方法用下列抽样分布来估计 TMLE 的抽样分布：从经验概率分布 $P_n$ 抽样，固定初始估计量，只重新计算 TMLE 算法中的 TMLE 更新步骤，取如此得到的 TMLE 的抽样分布。由此可知，只要 TMLE 本身渐近有效，这就是一种相合的自助法。当然，这类自助法无法捕捉因估计干扰参数而产生的二阶项。
